\documentclass[11pt]{article}
\usepackage[a4paper,margin=1in]{geometry}
\usepackage{amsmath,amssymb,amsthm,hyperref}
\newtheorem{theorem}{Theorem}
\newcommand{\Gr}{\operatorname{Gr}}
\newcommand{\tr}{\operatorname{tr}}
\newcommand{\diag}{\operatorname{diag}}
\title{Gaussian counterexamples to benign non-convexity\newline
in dimensional logarithmic Sobolev optimization}
\author{Principia Math In Progress}
\date{22 September 2026}
\begin{document}
\maketitle
\begin{abstract}
The dimensional logarithmic Sobolev majorant can have a strict local
minimum which is not global, even for a centered Gaussian target in
dimension two. We also give a separate Gaussian perturbation argument
for the nonequivalent sign convention printed in equation (18) of the
source. Thus the benign non-convexity conjecture fails under both readings.
\end{abstract}

\section*{Task 1: scope and the intended objective}
Conjecture 2.6 of Li, Cui, Li, Marzouk and Zahm \cite{source} asserts
benign non-convexity on $\Gr(d,r)$, meaning that every local minimum is
global. We take the permitted reference measure $\mu=N(0,I_d)$.
For a target $\pi$, let $M=\mathbb E_\pi[XX^{\mathsf T}]$ and
$H=\mathbb E_\pi[\nabla\log\pi\,\nabla\log\pi^{\mathsf T}]$.
The source's equation (20) gives the majorant
\[
 J(U)=\frac12\left\{\tr(U_\perp^{\mathsf T}MU_\perp)
 -(d-r)+\log\det(U_\perp^{\mathsf T}HU_\perp)\right\},       \tag{1}
\]
where $[U,U_\perp]$ is orthogonal. This is also the expression derived
in the source's Appendix D.2. In particular, its Proposition 4
identifies (1) with the actual approximation KL divergence for Gaussian
targets. We verify that identity directly in our example below.

There is a source sign inconsistency. Appendix B correctly states
\[
 \det(U_\perp^{\mathsf T}HU_\perp)
 =\det(H)\det(U^{\mathsf T}H^{-1}U).
\]
Consequently, the coefficient of $\log\det(U^{\mathsf T}H^{-1}U)$ in
(1) is $+1/2$. The displayed equation (18) instead has coefficient
$-1/2$. We first disprove the intended majorant landscape using (1),
and then separately disprove the landscape assertion for (18) as printed.
The resolution therefore does not rest on detecting the sign error.

\section*{A strict nonglobal minimum for the intended majorant}
\begin{theorem}
For $d=2$, $r=1$, $\mu=N(0,I_2)$ and
\[
 \pi=N(0,C),\qquad C=\diag(1/2,2),
\]
both coordinate lines are strict local minima of $J$ on $\Gr(2,1)$.
The line spanned by $e_1$ is not a global minimum.
\end{theorem}
\begin{proof}
Here $M=C$, $H=C^{-1}=\diag(2,1/2)$, and $\det C=1$.
For $u(\theta)=(\cos\theta,\sin\theta)^{\mathsf T}$ put
\[
 z(\theta)=u(\theta)^{\mathsf T}Cu(\theta)
 =\frac12+\frac32\sin^2\theta.
\]
Its orthogonal complement is $v(\theta)=(-\sin\theta,\cos\theta)^{\mathsf T}$.
Then $v^{\mathsf T}Cv=5/2-z$ and $v^{\mathsf T}Hv=z$, so (1) becomes
\[
 J(\theta)=\frac12\left(\frac32-z(\theta)+\log z(\theta)\right). \tag{2}
\]
The angle is a local coordinate on $\Gr(2,1)$, with period $\pi$.
At $\theta=0$ and $\theta=\pi/2$, the first derivative vanishes, and
direct differentiation gives
\[
 J''(0)=\frac32>0,\qquad J''(\pi/2)=\frac34>0.
\]
Thus both are strict local minima in the manifold itself. Their values are
\[
 J(0)=\frac{1-\log2}{2},\qquad
 J(\pi/2)=\frac{-1/2+\log2}{2}.
\]
The difference is $3/4-\log2>0$. For an elementary exact sign proof,
the graph of the strictly convex function $1/x$ lies strictly below its
secant on $(1,2)$; hence
\[
 \log2=\int_1^2\frac{dx}{x}
 <\int_1^2\frac{3-x}{2}\,dx=\frac34.
\]
The other coordinate line therefore has smaller value than the strict
local minimum at $e_1$, disproving benign non-convexity.
\end{proof}

These are valid probability measures under the source assumptions.
Their density ratio is the smooth positive function
$d\pi/d\mu=\exp(-x_1^2/2+x_2^2/4)$, normalized under $\mu$.
All second moments, Fisher information and KL quantities used above are finite.
The source assumes a smooth positive ratio, not a bounded ratio or a
target covariance dominated by $I_2$.

To check the actual KL interpretation independently, the source's optimal
profile preserves the $u$-marginal of $\pi$ and keeps a standard Gaussian
in the complementary direction. Its covariance is
\[
 D=I_2+(z-1)uu^{\mathsf T},\qquad
 D^{-1}=I_2+(z^{-1}-1)uu^{\mathsf T},\qquad \det D=z.
\]
The centered Gaussian KL formula yields
\[
 D_{\rm KL}(N(0,C)\Vert N(0,D))
 =\tfrac12\{\log z-\log\det C-2+\tr(D^{-1}C)\}
 =\tfrac12(3/2-z+\log z).
\]
Thus (2) is the intended approximation loss itself, independently of
the inconsistent display (18).

\section*{The literal equation (18) also has nonglobal local minima}
For completeness, write the function actually printed in (18) as
\[
 J_{\rm lit}(u)=\frac{\tr M-2+\log\det H}{2}
 -\frac{u^{\mathsf T}Mu-1+\log(u^{\mathsf T}H^{-1}u)}2.       \tag{3}
\]
This is again well defined on $\Gr(2,1)$.
\begin{theorem}
Let $C=\diag(1/2,2)$ and $m_\delta=(\sqrt3,\delta)^{\mathsf T}$.
There exists $\delta_0>0$ such that, for every $0<\delta<\delta_0$,
the target $\pi_\delta=N(m_\delta,C)$ gives two strict local minima of
(3) with different values. In particular, (3) is not benignly non-convex.
\end{theorem}
\begin{proof}
For these targets, $H=C^{-1}$ and $M=C+m_\delta m_\delta^{\mathsf T}$,
so all matrices in (3) arise from actual Gaussian probability measures.
Use the local chart
\[
 u(t)=\frac{(1,t)^{\mathsf T}}{\sqrt{1+t^2}},\qquad
 z(t)=\frac{1/2+2t^2}{1+t^2}.
\]
Then (3), as a smooth function of $(t,\delta)$, is
\[
 F(t,\delta)=\frac{9/2+\delta^2}{2}
 -\frac12\left\{
 \frac{1/2+2t^2+(\sqrt3+\delta t)^2}{1+t^2}+\log z(t)\right\}.
                                                               \tag{4}
\]
At $\delta=0$, set $w=(1+t^2)^{-1}$. Up to its additive constant,
(4) equals
\[
 -\tfrac12\{2+\tfrac32 w+\log(2-\tfrac32w)\}.
\]
Its derivative in $w$ vanishes at $w=2/3$, where its second derivative
is $9/8>0$. Accordingly, the two distinct lines with
$t_+=1/\sqrt2$ and $t_-=-1/\sqrt2$ satisfy
\[
 F_t(t_\pm,0)=0,\qquad F_{tt}(t_\pm,0)=\frac49>0.
\]
The implicit function theorem applied separately to $F_t=0$ provides
smooth critical points $t_\pm(\delta)$ for $|\delta|$ sufficiently small.
They remain distinct strict local minima by continuity of $F_{tt}$.
Let $v_\pm(\delta)=F(t_\pm(\delta),\delta)$ be their values.
Stationarity eliminates the derivatives of $t_\pm(\delta)$, giving
\[
 v_\pm'(0)=F_\delta(t_\pm,0)
 =-\frac{\sqrt3\,t_\pm}{1+t_\pm^2}
 =\mp\frac{\sqrt6}{3}.
\]
At zero the two values are equal, whereas
$(v_+-v_-)'(0)=-2\sqrt6/3<0$. It follows from the definition of the
derivative that $v_+(\delta)<v_-(\delta)$ for every sufficiently small
positive $\delta$. Thus the second local minimum is nonglobal, as claimed.
\end{proof}

Both the intended majorant and the literal printed function therefore
fail the original landscape assertion. These are two readings of the
same target, and receive only one problem credit.

\section*{Proposed follow-up Tasks 2--5}
These are proposed tasks, not additional completed results.
\begin{enumerate}
\setcounter{enumi}{1}
\item Classify all local minima of the intended objective on
$\Gr(d,r)$ for centered Gaussian targets with diagonal covariance.
\item Determine whether covariance domination $C\preceq I_d$ suffices
for benign non-convexity for every Gaussian target, including nonzero means.
\item Give an explicit numerical threshold $\delta_0$ in the second
theorem with certified lower bounds on both local Hessians.
\item Characterize pairs of second-moment and Fisher-information
matrices arising from probability measures for which (1) has no nonglobal
local minima for any rank $r$.
\end{enumerate}

\begin{thebibliography}{9}
\bibitem{source} M. T. C. Li, T. Cui, F. Li, Y. Marzouk and O. Zahm,
\emph{Sharp detection of low-dimensional structure in probability
measures via dimensional logarithmic Sobolev inequalities},
Information and Inference: A Journal of the IMA 14(3) (2025), iaaf021,
\href{https://doi.org/10.1093/imaiai/iaaf021}{doi:10.1093/imaiai/iaaf021}.
Complete author's final manuscript read from
\href{https://dspace.mit.edu/handle/1721.1/167023}{DSpace@MIT};
byte-identical to \href{https://arxiv.org/abs/2406.13036v3}{arXiv:2406.13036v3}.
\end{thebibliography}
\end{document}
